\section{Introduction}\label{sec:introduction}
An exact conic reformulation represents a convex body as the projection of
an affine slice of a product of cones. It can change the number of factors,
their dimensions, the rank of a supporting certificate, and the barrier
on the affine domain. It can also change which data a Newton step needs
and which outputs are inexpensive to produce. These resources are related,
but none determines all the others. The purpose of this paper is to give
precise comparisons, sharp formulas in specified models, and explicit
examples that separate the models.

Euclidean balls and their products provide the main setting. Their slack
has a simple differential structure, yet their lifts already exhibit
certificate sharing, singular boundary choices, nontrivial projection
fibers, and different barrier metrics. We also treat norm balls, matrix
balls, balance slices, power and entropy cones, and exact constraint
compilers. Throughout, an exact formula is accompanied by the class over
which it is optimized. A lower bound for a chosen representation is not
automatically a lower bound for every representation of its projection.

\subsection{Principal results and their scope}
The basic geometric calculation differentiates a nonnegative cone pairing
at a primal--polar contact. A proper cone of real dimension $m$ supplies
at most $\max\{m-2,0\}$ independent mixed-curvature directions. In a
simple Euclidean Jordan algebra, complementary primal and dual ranks
$p,q$ sharpen this count to $apq$, where $a$ is the Peirce constant
(Theorem~\ref{thm:curvature}). Definable local selections suffice;
no globally smooth selection is needed. For a ball of dimension $s\geq2$
and an integer cone-dimension cap $d\geq3$, the exact minimum factor count and
total cone dimension over exact definable lifts are consequently
\[
 k=\left\lceil\frac{s-1}{d-2}\right\rceil,
 \qquad M=s-1+2k.
\]
Classical norm trees attain both bounds
(Corollary~\ref{cor:dimension-resources}).

Our central rank invariant minimizes over the \emph{entire fiber of
genuine certificates for a fixed support}. A genuine certificate represents
the support slack on the whole affine slice, rather than only on selected
boundary points. Fix a finite repeatable dictionary $\mathscr K$ of simple
symmetric cones and put $B=\max_{K\in\mathscr K}a(K)(r(K)-1)>0$.
For every $s\geq2$, Theorem~\ref{thm:rank-frontier} proves
\begin{equation}\label{eq:intro-rank}
 \inf_{\mathcal L}\ \max_{v\in\Sph^{s-1}}
 \min_{Y\in\mathcal D_{\mathcal L}(v)}\sum_i\jrank Y_i
 =\left\lceil\frac{s-1}{B}\right\rceil.
\end{equation}
The lift is fixed before the support is chosen. For each fixed lift the
lower bound holds for every certificate on a dense open set of supports
with surface-null complement. A Peirce-perspective construction attains
the bound in every simple family, including the Albert cone. It also
shows why choosing the formulation after the objective changes the answer.
Sequential compression and private curvature describe which rank costs
add for product sources (Section~\ref{sec:products}).

Global selection is a different requirement. If primal boundary points
and dual certificates must be chosen $C^1$ over the entire support
manifold, saturation of a local rank bound forces a global map into a
product of support orbits or cone-base boundaries. Covering and cohomology
arguments then impose additional resource costs. We prove exact Lorentz
frontiers and selected symmetric-cone rank frontiers, together with
product saturation and whole-fiber rigidity theorems
(Sections~\ref{sec:regularity}--\ref{sec:full-fiber-rigidity}). The
whole-fiber results derive regular structure from a stronger fiber
hypothesis; they do not assume that arbitrary lifts have smooth selectors.

The barrier analysis distinguishes three optimization problems: the
standard product barrier in the ambient cone, that same barrier after
affine restriction, and the best barrier of any form on a fixed affine
domain. Boundary nullity and recession give lower tests for the restricted
standard parameter. These tests yield exact minima under global selection,
several selection-free exact cases, and intrinsic optima for grouped,
packed, and balance slices. Norm trees have a separate root-incidence law.
The results do not constitute a complete classification. For the remaining
bounded narrow-cap cases in \eqref{eq:unsettled-range}, the family is a
fixed-field Hermitian order cap or a Lorentz dimension cap, and $B$ is
that cap's primitive capacity. With $s=qB+1$ and $q\geq B+2$, the proved interval is
$q\leq\inf_{\mathcal L}\nustd\leq q+1$.
The nonsymmetric estimates in Section~\ref{sec:nonsymmetric-barriers}
likewise include parameter intervals, not general formulas for their optima.

Movement is always measured in a specified Hessian metric. Exposing
certificates yield path-independent target-distance bounds, and the rank
frontier makes these bounds depend on the permitted lift dictionary.
Primal--dual movement for arbitrary logarithmically homogeneous barriers
has a different endpoint and metric contract from primal movement in a
displayed barrier. An explicit refinement determines the shortest-distance
coefficient for every fixed weighted norm tree and fixed unit leaf objective
$c$. If $\omega_v\geq1$ is the barrier weight of internal node $v$, call
$v$ active when its subtree contains nonzero objective support. For the
target $1-c^Tw\leq\epsilon$, Theorem~\ref{thm:tree-exact-distance}
gives, from any fixed interior reference,
\begin{equation}\label{eq:intro-tree}
 d(\epsilon)=\sqrt{K_c}\log(1/\epsilon)
              +O(\log\log(1/\epsilon)),\qquad
 K_c=\sum_{v\text{ active}}\omega_v
                  +\frac12\sum_{v\text{ inactive}}\omega_v.
\end{equation}
The corresponding central route has coefficient
$\sqrt{\sum_v\omega_v}$. Thus inactive subtrees contribute to unavoidable
distance even though their exposing dual blocks vanish, and they also
produce a precisely quantified excess in central-path length. The tree,
weights, and objective are fixed before $\epsilon$ tends to zero.

For packed PSD representations, exact auxiliary elimination identifies a
common marginal Newton system at the fiber center. Away from that center,
the source allocation equations improve a generic squared comparison to
a sharp first-power transfer. At the same projected point, if the feasible
residuals and their exact diagonal center obey
$\mu D^0\preceq D\preceq L D^0$, where $0<\mu\leq1\leq L$,
Theorem~\ref{newt:linear} proves
\begin{equation}\label{eq:intro-transfer}
 L^{-1}H_0\preceq H_D\preceq\mu^{-1}H_0,
 \qquad\kappa(H_0^{-1/2}H_DH_0^{-1/2})\leq L/\mu.
\end{equation}
A two-source feasible example attains the condition factor. Exact
allocations are essential: approximate radial metadata needs its own
allocation error factor. The matrix comparison alone neither identifies
the reduced Newton right-hand side nor constructs a cheap preconditioner.

Finally, we specify when mathematical resources become computational
costs. Exact resource ledgers distinguish ambient coordinates, free
coordinates, and restricted parameters. Sparse elimination gives
classical comparisons under matched access and full-output contracts.
Query examples separate normalized projected states from readable values,
full vectors, exact reusable compilers, and fresh scale maintenance.
Some of these states need one query while classical output requires
linearly many queries. Our rank-dependent work bound retains the same
exposing certificate in the curvature and movement estimates and assumes
an explicit fresh charge per round (Theorem~\ref{work:total}). A one-shot query bound
cannot simply be multiplied by a number of geometric checkpoints.

\begin{table}[p]
\centering\small
\caption{Main results by model. Each cited statement gives its full
hypotheses. ``Exact'' refers only to the indicated optimization class.}
\label{tab:models}
\renewcommand{\arraystretch}{1.18}
\begin{tabular}{@{}>{\raggedright\arraybackslash}p{0.27\textwidth}>{\raggedright\arraybackslash}p{0.68\textwidth}@{}}
\toprule
Model and resource & Conclusion and scope\\
\midrule
Exact definable lifts; factor count and cone dimension
& Positive boundary curvature gives the dimension-minus-two bound.
Ball norm trees attain the capped count and dimension minima
(Theorem~\ref{thm:curvature}, Corollary~\ref{cor:dimension-resources}).\\
Symmetric-cone lifts; minimum certificate rank
& Entire-fiber minimization, then worst support, then lift minimization
gives $\lceil(s-1)/B\rceil$ without global selection
(Theorem~\ref{thm:rank-frontier}).\\
Globally $C^1$ full-slack selections; rank and factor resources
& Exact Lorentz and selected symmetric-cone frontiers; saturation forces
support-orbit or cone-base topology. This is a stronger model than
arbitrary affine lifts (Theorems~\ref{thm:smooth-lorentz},
\ref{thm:smooth-rank}, \ref{thm:smooth-product-rank}).\\
Restricted standard barrier; arbitrary barrier on a fixed domain
& Boundary/recession lower tests, delimited exact lift minima, and
concrete intrinsic optima. The bounded interval
\eqref{eq:unsettled-range} remains unresolved
(Sections~\ref{sec:restricted}, \ref{sec:concrete-barriers},
\ref{sec:balance-slices}).\\
Fixed barrier metric; distance to an accurate target
& Exposed-minor bounds and exact weighted tree coefficient $\sqrt{K_c}$.
The separate all-LH result concerns a central start and feasible
primal--dual gap endpoint
(Theorems~\ref{thm:movement-minor}, \ref{thm:tree-exact-distance},
\ref{thm:pd-gap-movement}).\\
Feasible PSD packing; projected Hessian quotient
& Exact centered reduction and sharp relative condition factor $L/\mu$
at the same projected point, with the source allocations fixed
(Proposition~\ref{newt:center}, Theorem~\ref{newt:linear}).\\
Specified access, output, and fresh-work contracts
& Conditional work composition, sparse full-output comparisons,
state/readout separations, and exact compiler acquisition costs
(Sections~\ref{work:section}--\ref{active:section}). Fresh adversarial
updates in Theorem~\ref{newt:dynamic-scale} are not a fixed program's
Newton trajectory.\\
\bottomrule
\end{tabular}
\end{table}

\subsection{Relation to the literature}
Conic lifts and slack factorizations are the starting framework of
Gouveia, Parrilo, and Thomas~\cite[Theorem~2.4]{GPT2013}; the survey
\citet{FGPPST2022Survey} places their use among geometric and algebraic
methods for concise convex descriptions. Fixed-size PSD factor counts
\citep{FP2013blocks}, PSD rank and algebraic-boundary bounds
\citep{FGPRT2015,FSED2018}, and semidefinite extension degree
\citep{Averkov2019} measure related but different resources. Extension
degree minimizes the largest permitted PSD block while allowing finite
products. Fawzi's SOC obstruction and Saunderson's short-face-chain
theorem use neighborliness to exclude finite product lifts in their
respective settings~\citep{Fawzi2019SOC,Saunderson2020}. Our local
curvature bounds quantify resources when a smooth body is representable;
our certificate invariant also minimizes within each fixed support fiber.
Kummer's ball bounds concern direct spectrahedral descriptions
\citep{Kummer2016}, which must be distinguished from projections of
products of cones.

Positive curvature itself does not obstruct SOC representability.
Scheiderer~\cite[Theorem~1.2]{Scheiderer2025SOC} proves that compact
convex semialgebraic sets with Nash-smooth boundary and everywhere strictly
positive curvature have SOC lifts. That existence theorem allows an
arbitrary finite number of order-two PSD blocks. It does not specify the
factor-count frontier, support-fiber minimax, or global-selection
frontiers studied here. This distinction is particularly relevant because
both approaches apply to positively curved boundaries.

The Jordan classification, Peirce spaces, quadratic representations, and
support manifolds are classical~\citep{FK1994,Nomura1994,GS2010}.
Symmetric-cone factorization algorithms provide additional context
\citep{SV2024}. Recursive small-Lorentz-cone norm representations are
also classical~\cite[Section~2]{BTN2001}. We use standard covering,
bundle, and cohomology tools~\citep{Hatcher2002,Steenrod1951}; their
application here depends on the precise global contact hypotheses.
The claimed geometric contributions are the stated quantitative
inequalities, fiber quantifiers, rigidity consequences, and attaining
constructions, not these underlying tools.

Self-concordant barrier theory supplies the analytic framework
\citep{NN1994}. Ambient optimality on symmetric cones is established
in the work of G\"uler--Tun\c{c}el and Cardoso--Vieira
\citep{GT1998,CardosoVieira2006}; self-scaled barriers have a classified
form~\citep{HauserGuler2002}. Conic extension, affine restriction, and
projective transformation of barriers belong to an existing theory
\citep{Hildebrand2022Projective}. Accordingly, the distinction between
ambient and restricted parameters is not itself a claim of novelty.
Our contribution is to compute the particular restricted or intrinsic
quantities proved here and connect them to exact lift resources.
For nonsymmetric cones we use established lower-bound and barrier
constructions~\citep{Hildebrand2013,Chares2009,FS2023}. The universal
barrier's parameter bound~\citep{LeeYue2021} and practical natural-cone
methods~\citep{CKV2022natural} also make clear why existence of a small
parameter should be separated from derivative evaluation cost.

Nesterov--Todd and Nesterov--Nemirovski develop the barrier geometry and
the comparison of central paths with Riemannian distances
\citep{NT2002,NesterovNemirovski2008}. The unpublished companion
\emph{The Cost of Following the Central Path}~\citep{CentralPathCompanion}
already contains the weighted exposed-minor estimate, several
distribution-sensitive movement bounds, and standard-barrier calculations
for grouped, packed, and norm-tree formulations. We reproduce the proofs
needed here, so no theorem depends on access to that companion. The
formulation-specific developments here include composition with the
certificate-rank frontier, the general affine-PSD exposure profile, and
the exact active/inactive weighted-tree coefficient
\eqref{eq:intro-tree}. We do not claim the general distinction between
central-path length and shortest distance as a new observation.

The computational arguments likewise use established ingredients:
search, oracle interrogation, approximate counting, and information
bounds~\citep{BBBV1997,vanDam1998,NayakWu1999,Nayak1999,BHMT2002};
sparse cone augmentation~\citep{DomahidiChuBoyd2013}; and elimination
from a supplied tree decomposition~\citep{FurerHoppenTrevisan2025}.
Our statements attach these tools to specified conic representations and
outputs. Quantum interior-point methods and their end-to-end resource
analyses~\citep{AugustinoEtAl2023,DalzellEtAl2023} are relevant algorithmic
comparators. More recent iterative refinement improves the high-precision
behavior of quantum interior-point methods
\citep{MohammadisiahroudiEtAl2026}; we make no blanket claim that all such
methods retain earlier inverse-accuracy dependencies. Our exact-field
full-output comparisons and the feasible quotient theorem
\eqref{eq:intro-transfer} have their own explicit access, feasibility,
and output assumptions. In particular, the first-power improvement uses
elementary positive trace inequalities and the allocation equalities;
it is not a new general matrix-order principle.

The specific results we put forward as new are the selection-free
entire-fiber minimax \eqref{eq:intro-rank}, the active/inactive weighted-tree
coefficient \eqref{eq:intro-tree}, and the sharp feasible allocation-aware
quotient transfer \eqref{eq:intro-transfer}, in the precise models stated
above. These claims concern the formulas and their quantifiers, not priority
for the general methods used to prove them. The comparisons above identify
the scope of the checked primary results; absence from a literature search
is not evidence of priority.

\subsection{Conventions and reading guide}
Table~\ref{tab:notation} collects the recurring notation. Symbols such as
$p,q,L$ have local uses in some construction and counting sections; those
uses are defined where introduced. A theorem's stated domain, metric, and
order of optimization always take precedence over an informal resource
name.

\begin{table}[htbp]
\centering\small
\caption{Recurring notation and conventions.}\label{tab:notation}
\renewcommand{\arraystretch}{1.12}
\begin{tabular}{@{}>{\raggedright\arraybackslash}p{0.27\textwidth}>{\raggedright\arraybackslash}p{0.68\textwidth}@{}}
\toprule
Notation & Meaning\\\midrule
$\mathcal L=(K,A,\pi)$ & Exact affine lift; $K=\prod_iK_i$.
Relative Slater reduction is made before computing ranks and barriers
(Lemma~\ref{lem:reduction}).\\
$\mathcal D_{\mathcal L}(v)$ & Positive certificates representing
$1-\langle\pi(X),v\rangle$ for every $X\in A$,
including infeasible slice points.\\
$s$, $s_a$; $m_i$ & Projected ball dimensions; real cone-space dimensions.
A Lorentz cone $Q_m$ has total dimension $m$.\\
$r_i,a_i$; $p_i,q_i$ & Jordan rank and Peirce constant of a simple factor;
complementary primal and dual contact ranks. $\jrank$ denotes Jordan
rank; real PSD Jordan rank is ordinary matrix rank.\\
$B$; $a_i\lfloor r_i^2/4\rfloor$ & Maximum primitive capacity per dual
rank in a dictionary; maximum mixed-curvature capacity of one factor.
These capacities are different.\\
$\nustd$ & Least parameter of the specified standard log-determinant
barrier on its effective affine slice; see
\eqref{eq:exact-gradient-parameter}.\\
$\vartheta_{\mathrm{opt}}(\Omega)$;
$\nu_{\mathrm{LH}}(C)$ & Infimum over all self-concordant barriers on a
fixed affine domain; infimum restricted to logarithmically homogeneous
barriers on a fixed cone, respectively.\\
$d_F$, $\epsilon$ & Distance in the Hessian metric of the displayed
barrier; objective-gap or output accuracy as locally specified.
Distances to targets minimize over paths in the stated metric domain,
with the endpoint constraints of the relevant theorem.\\
PSD order; asymptotic notation & $H\preceq G$ means $G-H$ is positive
semidefinite. Constants and fixed-instance versus joint dimension/accuracy
limits are stated in each result.\\
\bottomrule
\end{tabular}
\end{table}

Part~\ref{part:local} develops local curvature and certificate fibers.
Part~\ref{part:regularity} treats global selection, restricted barriers,
and saturation. Part~\ref{part:families} gives formulation families and
nonsymmetric comparisons. Part~\ref{part:compilation} addresses controlled
approximation and exact compilation. Part~\ref{part:movement} connects
certificates to metric distance, culminating in the weighted tree theorem.
Part~\ref{part:computation} gives resource ledgers and computational
contracts, including the sharp PSD quotient comparison. The appendices
retain topological refinements, restricted-kernel distinctions, robust
approximation, and conditioning counterexamples. Each part supplies the
proofs needed for its formulation claims; the contents and model table
allow readers to follow a particular resource without treating every
result as a theorem about the same class of lifts.
